\fancyhead[L]{\small THEOREMATIS DE RESOLUBILITATE}
\SourceLeaf{066}
\begin{center}
\textbf{THEOREMATIS DE RESOLUBILITATE}\\[8pt]
\textbf{FUNCTIONUM ALGEBRAICARUM INTEGRARUM}\\[8pt]
\textbf{IN FACTORES REALES}\\[8pt]
DEMONSTRATIO TERTIA.\\[8pt]
SUPPLEMENTUM COMMENTATIONIS PRAECEDENTIS.\\[12pt]
\rule{3cm}{0.3pt}
\end{center}
Postquam commentatio praecedens typis iam expressa esset, iteratae de eodem argumento meditationes ad novam theorematis demonstrationem perduxerunt,
quae perinde quidem ac praecedens pure analytica est, sed principiis prorsus diversis innititur, et respectu simplicitatis illi longissime praeferenda videtur. Huic
itaque \textit{tertiae} demonstrationi pagellae sequentes dicatae sunto.

\ThirdArticle{1}
Proposita sit functio indeterminatae $x$ haecce:
\[
X=x^m+Ax^{m-1}+Bx^{m-2}+Cx^{m-3}+\text{etc.}+Lx+M
\]
in qua coëfficientes $A,B,C$ etc. sunt quantitates reales determinatae. Sint $r,\varphi$
aliae indeterminatae, statuamusque
\[
\begin{aligned}
&r^m\cos m\varphi+Ar^{m-1}\cos(m-1)\varphi\\
&\qquad+Br^{m-2}\cos(m-2)\varphi\\
&\qquad+Cr^{m-3}\cos(m-3)\varphi+\text{etc.}+Lr\cos\varphi+M=t\\
&r^m\sin m\varphi+Ar^{m-1}\sin(m-1)\varphi\\
&\qquad+Br^{m-2}\sin(m-2)\varphi\\
&\qquad+Cr^{m-3}\sin(m-3)\varphi+\text{etc.}+Lr\sin\varphi=u
\end{aligned}
\]
\SourceLeaf{067}
\[
\begin{aligned}
&mr^m\cos m\varphi+(m-1)Ar^{m-1}\cos(m-1)\varphi\\
&\qquad+(m-2)Br^{m-2}\cos(m-2)\varphi\\
&\qquad+(m-3)Cr^{m-3}\cos(m-3)\varphi+\text{etc.}+Lr\cos\varphi=t'\\
&mr^m\sin m\varphi+(m-1)Ar^{m-1}\sin(m-1)\varphi\\
&\qquad+(m-2)Br^{m-2}\sin(m-2)\varphi\\
&\qquad+(m-3)Cr^{m-3}\sin(m-3)\varphi+\text{etc.}+Lr\sin\varphi=u'\\[6pt]
&mmr^m\cos m\varphi+(m-1)^2Ar^{m-1}\cos(m-1)\varphi\\
&\qquad+(m-2)^2Br^{m-2}\cos(m-2)\varphi\\
&\qquad+(m-3)^2Cr^{m-3}\cos(m-3)\varphi+\text{etc.}+Lr\cos\varphi=t''\\
&mmr^m\sin m\varphi+(m-1)^2Ar^{m-1}\sin(m-1)\varphi\\
&\qquad+(m-2)^2Br^{m-2}\sin(m-2)\varphi\\
&\qquad+(m-3)^2Cr^{m-3}\sin(m-3)\varphi+\text{etc.}+Lr\sin\varphi=u''
\end{aligned}
\]
\[
\frac{(tt+uu)(tt''+uu'')+(tu'-ut')^2-(tt'+uu')^2}{r(tt+uu)^2}=y
\]
Factorem $r$ manifesto e denominatore formulae ultimae tollere licet, quum
$t',u',t'',u''$ per illum sint divisibiles. Denique sit $R$ quantitas positiva determinata, arbitraria quidem, attamen maior maxima quantitatum
\[
mA\BareRoot{2},\quad\BareRoot{(mB\BareRoot{2})},\quad\BareRoot[3]{(mC\BareRoot{2})},\quad\BareRoot[4]{(mD\BareRoot{2})}\quad\text{etc.}
\]
abstrahendo a signis quantitatum $A,B,C$ etc., i. e. mutatis negativis, si quae
adsint, in positivas. His ita praeparatis, dico, $tt'+uu'$ certo nancisci valorem
positivum, si statuatur $r=R$, quicunque valor (realis) ipsi $\varphi$ tribuatur.

\textit{Demonstratio.} Statuamus
\[
\begin{aligned}
&R^m\cos45^\circ+AR^{m-1}\cos(45^\circ+\varphi)\\
&\qquad+BR^{m-2}\cos(45^\circ+2\varphi)\\
&\quad+CR^{m-3}\cos(45^\circ+3\varphi)+\text{etc.}\\
&\qquad+LR\cos(45^\circ+(m-1)\varphi)+M\cos(45^\circ+m\varphi)=T\\
&R^m\sin45^\circ+AR^{m-1}\sin(45^\circ+\varphi)\\
&\qquad+BR^{m-2}\sin(45^\circ+2\varphi)\\
&\quad+CR^{m-3}\sin(45^\circ+3\varphi)+\text{etc.}\\
&\qquad+LR\sin(45^\circ+(m-1)\varphi)+M\sin(45^\circ+m\varphi)=U\\
&mR^m\cos45^\circ+(m-1)AR^{m-1}\cos(45^\circ+\varphi)\\
&\qquad+(m-2)BR^{m-2}\cos(45^\circ+2\varphi)\\
&\quad+(m-3)CR^{m-3}\cos(45^\circ+3\varphi)+\text{etc.}\\
&\qquad+LR\cos(45^\circ+(m-1)\varphi)=T'\\
&mR^m\sin45^\circ+(m-1)AR^{m-1}\sin(45^\circ+\varphi)\\
&\qquad+(m-2)BR^{m-2}\sin(45^\circ+2\varphi)\\
&\quad+(m-3)CR^{m-3}\sin(45^\circ+3\varphi)+\text{etc.}\\
&\qquad+LR\sin(45^\circ+(m-1)\varphi)=U'
\end{aligned}
\]
patetque

I. $T$ compositam esse e partibus
\SourceLeaf{068}
\[
\begin{aligned}
&\frac{R^{m-1}}{m\BareRoot{2}}[R+mA\BareRoot{2}\, .\,\cos(45^\circ+\varphi)]\\
&+\frac{R^{m-2}}{m\BareRoot{2}}[RR+mB\BareRoot{2}\, .\,\cos(45^\circ+2\varphi)]\\
&+\frac{R^{m-3}}{m\BareRoot{2}}[R^3+mC\BareRoot{2}\, .\,\cos(45^\circ+3\varphi)]\\
&+\frac{R^{m-4}}{m\BareRoot{2}}[R^4+mD\BareRoot{2}\, .\,\cos(45^\circ+4\varphi)]\\
&+\text{etc.}
\end{aligned}
\]
quas singulas, pro valore quolibet determinato reali ipsius $\varphi$, positivas evadere
facile perspicitur: hinc $T$ necessario valorem positivum obtinet. Simili modo
probatur, etiam $U,T',U'$ fieri positivas, unde etiam $TT'+UU'$ necessario
fit quantitas positiva.

II. Pro $r=R$ functiones $t,u,t',u'$ resp. transeunt in
\[
\begin{gathered}
T\cos(45^\circ+m\varphi)+U\sin(45^\circ+m\varphi)\\
T\sin(45^\circ+m\varphi)-U\cos(45^\circ+m\varphi)\\
T'\cos(45^\circ+m\varphi)+U'\sin(45^\circ+m\varphi)\\
T'\sin(45^\circ+m\varphi)-U'\cos(45^\circ+m\varphi)
\end{gathered}
\]
uti evolutione facta facile probatur. Hinc vero valor functionis $tt'+uu'$, pro
$r=R$, derivatur $=TT'+UU'$, adeoque est quantitas positiva. Q. E. D.

Ceterum ex iisdem formulis colligimus valorem functionis $tt+uu$, pro $r=R$,
esse $TT+UU$, adeoque positivum, unde concludimus, pro nullo valore ipsius $r$,
singulis $mA\BareRoot{2},\BareRoot{(mB\BareRoot{2})},\BareRoot[3]{(mC\BareRoot{2})}$ etc. maiori,
simul fieri posse $t=0,u=0$.

\ThirdArticle{2}
\textsc{Theorema.} \textit{Intra limites $r=0$ et $r=R$, atque $\varphi=0$ et $\varphi=360^\circ$
certo exstant valores tales indeterminatarum $r,\varphi$, pro quibus fiat simul $t=0$ et $u=0$.}

\textit{Demonstratio.} Supponamus theorema non esse verum, patetque, valorem
ipsius $tt+uu$ pro cunctis valoribus indeterminatarum intra limites assignatos
fieri debere quantitatem positivam, et proin valorem ipsius $y$ semper finitum. Consideremus integrale duplex
\[
\iint y\,\mathrm dr\,\mathrm d\varphi
\]
ab $r=0$ usque ad $r=R$, atque a $\varphi=0$ usque ad $\varphi=360^\circ$ extensum,
quod igitur valorem finitum plene determinatum nanciscitur. Hic valor, quem
\SourceLeaf{069}
per $\Omega$ denotabimus, idem prodire debebit, sive integratio primo instituatur secundum $\varphi$ ac dein secundum $r$, sive ordine inverso. At habemus \textit{indefinite}, considerando $r$ tamquam constantem,
\[
\int y\,\mathrm d\varphi=\frac{tu'-ut'}{r(tt+uu)}
\]
uti per differentiationem secundum $\varphi$ facile confirmatur. Constans non adiicienda, siquidem integrale a $\varphi=0$ incipiendum supponamus, quoniam pro
$\varphi=0$ fit $\frac{tu'-ut'}{r(tt+uu)}=0$. Quare quum manifesto $\frac{tu'-ut'}{r(tt+uu)}$ etiam evanescat pro
$\varphi=360^\circ$, integrale $\int y\,\mathrm d\varphi$ a $\varphi=0$ usque ad $\varphi=360^\circ$ fit $=0$, manente $r$
indefinita. Hinc autem sequitur $\Omega=0$.

Perinde habemus indefinite, considerando $\varphi$ tamquam constantem,
\[
\int y\,\mathrm dr=\frac{tt'+uu'}{tt+uu}
\]
uti aeque facile per differentiationem secundum $r$ confirmatur: hic quoque constans non adiicienda, integrali ab $r=0$ incipiente. Quapropter integrale ab
$r=0$ usque ad $r=R$ extensum fit per ea, quae in art. praec. demonstrata sunt,
$=\frac{TT'+UU'}{TT+UU}$ adeoque per theorema art. praec. semper quantitas positiva pro quolibet valore reali ipsius $\varphi$. Hinc etiam $\Omega$, i. e. valor integralis
\[
\int\frac{TT'+UU'}{TT+UU}\,\mathrm d\varphi
\]
a $\varphi=0$ usque ad $\varphi=360^\circ$, necessario fit quantitas positiva\HistoricalNote{13}. Quod est absurdum, quoniam eandem quantitatem antea invenimus $=0$: suppositio itaque
consistere nequit, theorematisque veritas hinc evicta est.

\ThirdArticle{3}
Functio $X$ per substitutionem $x=r(\cos\varphi+\sin\varphi\, .\,\BareRoot{-1})$ transit in $t+u\BareRoot{-1}$,
nec non per substitutionem $x=r(\cos\varphi-\sin\varphi\, .\,\BareRoot{-1})$ in $t-u\BareRoot{-1}$. Quodsi
igitur pro valoribus determinatis ipsarum $r,\varphi$, puta pro $r=g,\varphi=G$, simul
provenit $t=0,u=0$ (quales valores exstare in art. praec. demonstratum est),
$X$ per utramque substitutionem
